\scanpage{13}
\mspage{11}
plus précisément appeler \uline{complexe de De Rham à pd virtuel} \inserted{\uline{sur $k$}} une \deleted{complexe de} $k$-bi-algèbre diff. \deleted{\unsure{k}}augmentée, ass. unit. alternée, à diff de bidegré $-1,+1$, \deleted{(ou $0,+1$ au choix)} à bidegrés $\geq 0$, avec p.d. sur l'idéal d'augm, compatibles avec le diff ($d(x^{(n)})=x^{(n-1)}dx$), et telle que les $H^{\bullet q}(C^{\bullet\bullet})[-q]$ \deleted{sont} des $q$-ombres \unsure{$\simeq$|$\sim$} (auquel cas on récupère à partir de $C^{\bullet\bullet}$ un élément de $D(k)$ avec structure multiplicative ass. unit. comm. $\ldots$), passer à une ``catégorie dérivée'' de ces complexes en inversant les flèches qui sont des \unsure{qu-isom|quasi-isom|q-isom.}, d'où par $C^{\bullet\bullet}_{\mathrm{DRpd}}(\unsure{\mathscr X|X|\mathcal X},k)$ un foncteur
\diagram{\begin{tikzcd}[ampersand replacement=\&]
(\mathrm{Hot})\arrow[r] \& (\mathrm{DRpd})\\[-2pt]
\text{types d'hom} \& \substack{\text{cat. dérivée}\\\text{des complexes}\\\mathrm{DR\ pd}}
\end{tikzcd}}
et on peut se demander si sa restriction aux espaces connexes et simplemt connexes avec des $H_{\deleted{\unsure{\ast|i}}}^{i}(X,\mathbb Z)$ de t.f. (cas $k=\mathbb Z$) \deleted{est} \inserted{induit} une équivalence avec les complexes de DR à pd sur $\mathbb Z$ tels que $H^0(C^{\bullet\bullet})\simeq k$, $H^1(C^{\bullet\bullet})=0$, les $H^i(C^{\bullet\bullet})$ de t.f. sur $\mathbb Z$ \ldots

% archive sheet number: 12
